% =============================================================================
% Project : THE_PHILOSOPHY_OF_VIABILITY
% Book    : Book I — Ontology of Viability
% Book ID : VS.B01
% Part    : Part I — Identity and Manifestation (VS.B01.P01)
% Chapter : Chapter 08 — Comparative Philosophical Examination
% Chapter ID : VS.B01.P01.C08
%
% File    : 4_ASSETS/diagrams/VS.B01.P01.C08.S03.SS04.P005.DG01.heraclitean_relations_of_opposites.tex
%
% Section / Chapter / Asset:
%   Diagram DG01 — Heraclitean Relations of Opposites
%
% Structural Identifier : VS.B01.P01.C08.S03.SS04.P005.DG01
% Asset Type : DG
%
% Purpose :
%   To distinguish four interpretively relevant forms of relation among
%   opposites without reducing them to one logical structure.
%
% Representational limit :
%   The categories are analytical guides for comparison, not a formal
%   taxonomy attributed directly to Heraclitus.
%
% Dependencies :
%   TikZ; libraries used: calc, positioning. All are loaded by the project
%   preamble; no per-file package or library loading.
%
% =============================================================================

\begin{figure}[htbp]
    \centering
    \hyphenpenalty=10000 \exhyphenpenalty=10000\relax  % no hyphenation inside diagram labels

    \begin{tikzpicture}[
        font=\small,
        node distance=1.3cm and 1.2cm,
        % --- base box ---
        box/.style={
            draw, rounded corners=3pt, align=center,
            text width=3.2cm, minimum height=1.4cm,
            inner sep=5pt, line width=0.7pt
        },
        % --- node roles ---
        relationbox/.style={box, fill=blue!8, draw=blue!55!black},
        centrebox/.style={box, text width=3.6cm, fill=blue!16,
                          draw=blue!55!black, font=\small\bfseries},
        % --- undirected link: no priority between the four relations ---
        vslink/.style={draw=black!60, line width=0.7pt}
    ]

    % ---------------------------------------------------------------
    % Central theme and four interpretive relations
    % ---------------------------------------------------------------
    \node[centrebox] (centre) {Relations among Opposites};

    \node[relationbox, above left=of centre]  (context) {
        \textbf{Contextual Contrast}\\[2pt] {\footnotesize B61}
    };
    \node[relationbox, above right=of centre] (reciprocal) {
        \textbf{Reciprocal Dependence}\\[2pt] {\footnotesize B10}
    };
    \node[relationbox, below left=of centre]  (alternation) {
        \textbf{Temporal Alternation}\\[2pt] {\footnotesize B88}
    };
    \node[relationbox, below right=of centre] (tension) {
        \textbf{Tension and Attunement}\\[2pt] {\footnotesize B51, B54}
    };

    \draw[vslink] (context)     -- (centre);
    \draw[vslink] (reciprocal)  -- (centre);
    \draw[vslink] (alternation) -- (centre);
    \draw[vslink] (tension)     -- (centre);

    % ---------------------------------------------------------------
    % Note (centred below the lower pair)
    % ---------------------------------------------------------------
    \coordinate (mid) at ($(alternation.south)!0.5!(tension.south)$);
    \node[below=0.7cm of mid, text width=11cm, align=center,
          font=\footnotesize\itshape] {
        These are distinguishable interpretive relations. The diagram
        does not assert that the opposed terms are identical, that every
        opposition has the same structure, or that the relations are
        exhaustive of Heraclitus' thought.
    };

    \end{tikzpicture}

    \caption{Four distinguishable ways in which Heraclitus' fragments
    present relations among opposites. The categories are
    analytical guides for comparison, not a formal taxonomy
    attributed directly to Heraclitus.}
    \label{fig:VS.B01.P01.C08.S03.SS04.P005.DG01}
\end{figure}

% =============================================================================
% END OF DIAGRAM DG01
% =============================================================================
